\subsection{An authenticated execution of the golden scalar action}
\label{sec:source-scalar-execution}

A small execution retains the field reads that realize the tensor action.
It uses the $q=5$ golden source coordinates
\[
 0,\quad 2\phi-3,\quad4\phi-6,\quad\phi-1,\quad3\phi-4,\quad1.
\]
The first five are the existing golden population; the final point is the
fixed Dirichlet boundary marker. The $4^3=64$ interior patches carry mutable
field registers. Their source addresses and action coefficients are protected
preparation data. Axes, boundary conditions, the action and model time remain
declared inputs.

Use the product masses $M$, tensor operator $\mathsf A$ and normalized
coordinates of \eqref{eq:common-scalar-action}. Thus
$\mathsf A$ is positive and self-adjoint in the $M$ norm, with
$\mathsf A\ge I$. Put
\[
 \Delta=1/\sqrt5,\qquad \tau=\Delta/7=(2\phi-1)/35,
 \qquad\tau^2=1/245.
\]
The discrete action and its exact Euler--Lagrange update are
\begin{align}
 \mathcal S_\tau&=\sum_j\left[
 \frac{\|u_{j+1}-u_j\|_M^2}{2\tau}
 -\frac\tau2\langle u_j,\mathsf A u_j\rangle_M\right],\notag\\
 u_{j+1}&=2u_j-u_{j-1}-\tau^2\mathsf A u_j.
 \label{eq:source-scalar-executed-update}
\end{align}
The step reads the previous local value and the current local value and
current nearest neighbours. Fixed boundary values contribute on-site terms.
Two alternating registers per patch retain immutable versions of the two
layers. A serialized event names every consumed resource, version, actual
writer and exact value, and commits its new value. The serial hash parent
is also retained; it is not inserted into the semantic field-read order.

The preparation has $u_0=0$ and
\[
 v_i=256\prod_{a=1}^3x_{ia}(1-x_{ia})\,
       \mathbf1_{x_{i1}<9/20},\qquad
 g_i=64\prod_{a=1}^3x_{ia}(1-x_{ia})\,
       \mathbf1_{x_{i1}>11/20}.
\]
The detector is $D(u)=\langle g,u\rangle_M$.
The initial memory $u_{-1}=-\tau v$ gives the centered initial velocity
$v$. It is integrator preparation, not a recorded exact solution at negative
time. This algebraic preparation differs from the nine-sine preparation in
Theorem~\ref{thm:source-common-scalar-response}; its numerical accuracy is
not transferred to this coarser population.

\begin{proposition}[Stable exact execution and preserved energy]
\label{prop:source-scalar-execution-energy}
For this action, $\tau^2\|\mathsf A\|_M<3<4$.
Every step of \eqref{eq:source-scalar-executed-update} preserves
\begin{align}
 \mathcal E_\tau(u_{j-1},u_j)
 &=\frac{\|u_j-u_{j-1}\|_M^2}{2\tau^2}
       +\frac12\langle u_j,\mathsf A u_{j-1}\rangle_M\notag\\
 &=\frac12\|\pi_j\|_M^2+
 \frac12\langle u_j,\mathsf A(I-\tau^2\mathsf A/4)u_j\rangle_M,
 \label{eq:source-scalar-execution-energy}
\end{align}
where $\pi_j=(u_j-u_{j-1})/\tau-\tau\mathsf A u_j/2$.
In particular the energy is positive, and
$\mathcal E_\tau\ge\|\pi_j\|_M^2/2+\|u_j\|_M^2/8$.
\end{proposition}
\begin{proof}
The smallest coordinate gap is $5-3\phi>7/50$. The row sums of the
absolute one-dimensional operator are bounded by $4/h_{\min}^2$.
Similarity to its symmetric mass-whitened matrix therefore gives
\[
 \|\mathsf A\|_M\le1+12/h_{\min}^2
 <30049/49<614,\qquad
 \tau^2\|\mathsf A\|_M<30049/12005<3.
\]
Subtract successive first-line energies and use the recurrence and
self-adjointness. The difference is zero. Completing the square gives
the second line, and $I-\tau^2\mathsf A/4> I/4$ proves coercivity.
The executable verifier checks this equality exactly at all twenty-two
retained states. It is a modified discrete energy; conservation of the
original continuous Hamiltonian is not asserted at nonzero step size.
\end{proof}

\paragraph{Full continuous-action comparison.}
The reference here is the continuous-time \emph{same finite} $64$-mode
operator, with the same $u_0=0,v_0=v$:
\[
 U(t)=\mathsf A^{-1/2}\sin(t\sqrt{\mathsf A})v.
\]
Its comparison to the execution includes every spectral mode. Let
\[
 P_{80}(t)=\sum_{k=0}^{80}
       \frac{(-1)^kt^{2k+1}}{(2k+1)!}\mathsf A^kv.
\]
The field data obey $0\le v_i\le4$, $0\le g_i\le1$ and
$\sum_iM_i\le1$, hence $\|v\|_M\le4$, $\|g\|_M\le1$.
Taylor's theorem for sine through degree $162$, followed by spectral
calculus and the preceding operator bound, gives for
$0\le t\le3/\sqrt5<3/2$
\begin{equation}
 \|U(t)-P_{80}(t)\|_M
 \le R:=\frac{4(3/2)^{163}614^{81}}{163!}<10^{-30}.
 \label{eq:source-scalar-execution-tail}
\end{equation}
The polynomial, the executed values and their squared $M$-norm difference
are exact elements of $\mathbb Q(\phi)$. Exact sign bisection encloses
their square roots. Adding $R$ gives a full state-error enclosure at every
executed time $t_j=j\tau$, $1\le j\le21$, and detector error at most
the same bound. The reference is not a truncated eigenvalue simulation.

The exact energy is approximately $0.276440884$. At $t=3/\sqrt5$, the
full state discretization error is enclosed between $0.009539009052$ and
$0.009539009055$. These are finite-operator comparisons at executed times,
not spatial continuum errors or errors for intermediate unstored times.
The independent verifier uses the $1,\sqrt5$ basis, assembles undirected
edge energies and evaluates the same full Taylor polynomial in reverse
order. It validates the scientific values as well as the commitments.

\paragraph{Canonical quantum compatibility.}
In mass-whitened canonical coordinates, the update is the Heisenberg map
of the symmetric unitary product
\[
 e^{-i\tau V/(2\hbar)}e^{-i\tau T/\hbar}
 e^{-i\tau V/(2\hbar)},\qquad
 T=\tfrac12P^TP,\quad V=\tfrac12X^TA_qX.
\]
Thus the exact classical impulse response also determines the central
commutator of the corresponding smeared canonical fields:
\[
 \frac{[\Phi(g)_j,\Phi(v/4)_0]}{i\hbar}
      =-\frac14D(u_j)\,I.
\]
This identity is state independent, initially on the common Schwartz
domain and then as its constant quadratic-form extension. It uses all
finite oscillator modes. No covariance or measurement probability is
computed. In particular no modified-Hamiltonian vacuum replaces the
original source-action vacuum. A probability comparison under these gates
would also have to propagate that original covariance.

\paragraph{Actual reads and the clock boundary.}
One trace stores $128$ preparation writes and $21\times64=1344$ update
writes. There are $8\cdot4^3-6\cdot4^2=416$ dynamic field reads per
update layer, or $8736$ in total. These counts include every previous-local
and current-layer field read; the immutable source and coefficient
commitments are supplied once as preparation inputs. Baseline and
intervention traces under ascending and descending site schedules have
identical corresponding field layers. Their full audit chains are retained
and differ with schedule.

Inclusive intervals in the actual field-read order, from the current-layer
seed at site $(1,1,1)$ to the same site after $7,14,21$ steps, contain
$144,576,1024$ events. They are not the intervals of the earlier all-neighbour
source-count law. For example the prepared source plane is at
$x=2\phi-3$ and the recorded target at $x=3\phi-4$ has a nonzero response
after seven substeps. Their point displacement squared is $2-\phi>1/5$,
exceeding the old radius squared at that coarse layer. This exposes a
finite point-event ancestry mismatch. It is not a continuum superluminality
claim: interpolation supports differ from source-point locations, and no
spatial continuum error bound is claimed at this population.

All substeps remain visible. Counting a chosen subset of their events does
not by itself preserve the causal intervals required for a count clock.
Relating this execution to that clock requires an explicit common event
or operational readout map with its intervention, order and volume errors.
The present trace supplies the inputs for that comparison; it does not
identify native time, primitive word costs or laboratory seconds.
